% ═══ Results: The Information Ceiling of a Feature Class ═══
\section{The Information Ceiling of a Feature Class}
\label{sec:res-ceiling}

Section~\ref{sec:res-transfer} reports that the universal circuit transfers to
held-out proteins at prec@$L/5 = 0.179$ against published DCA's $0.723$, and
Section~\ref{sec:res-reconstruction-negative} that this is insufficient for
template-free reconstruction.  Both are empirical statements about one method.
This section replaces them with a bound that holds for \emph{every} method
built on the same features, and thereby identifies where the shortfall lives.

The bound rests on an observation that is trivial and decisive: a circuit whose
inputs are a feature vector is a function of that feature vector, and therefore
cannot distinguish two candidate pairs that share it.

\begin{definition}[Cell predictor]
\label{def:cellpredictor}
Let $P$ be the candidate residue pairs of a protein, $y : P \to \{0,1\}$ the
native-contact labels, and $\varphi : P \to \mathcal{C}$ a feature map into a
finite set of \emph{cells}.  A \emph{cell predictor} is any
$\rho : \mathcal{C} \to \mathbb{R}$, inducing the score $\rho \circ \varphi$.
Pairs sharing a cell receive equal scores; their relative order is fixed by a
tie-break which carries no information about $y$ and is therefore modelled as
uniformly random.
\end{definition}

Every circuit in this paper is a cell predictor: a sum-of-products over
$\varphi$ is a $\{0,1\}$-valued function of the cell, and the smoothed
cell-rate scoring used in practice is a real-valued one.  The class of cell
predictors is strictly larger than the class of circuits, so a bound over it
binds the circuits a fortiori.

\begin{theorem}[Ceiling]
\label{thm:ceiling}
Write $n_c = |\varphi^{-1}(c)|$ and $r_c$ for the empirical contact rate in
cell $c$.  For any cell predictor $\rho$ and any $k$,
\[
\mathbb{E}\bigl[\operatorname{prec@}k(\rho \circ \varphi)\bigr]
\;\leq\; \frac{1}{k}\max\Bigl\{ \textstyle\sum_{c} t_c\, r_c \;:\;
0 \le t_c \le n_c,\ \sum_c t_c = k \Bigr\},
\]
the expectation taken over the tie-break, with the maximum attained by filling
cells in decreasing order of $r_c$.
\end{theorem}

\begin{proof}
Fix $\rho$.  Its induced ranking visits cells in decreasing $\rho$-order; write
$t_c$ for the number of pairs drawn from cell $c$ among the top $k$, so
$0 \le t_c \le n_c$ and $\sum_c t_c = k$.  Within a cell the tie-break is
uniform, so the $t_c$ selected pairs form a uniformly random $t_c$-subset of
$\varphi^{-1}(c)$ and the expected number of contacts drawn from $c$ is
$t_c r_c$ by linearity of expectation; hence
$\mathbb{E}[\operatorname{prec@}k] = k^{-1}\sum_c t_c r_c$.  Varying $\rho$
varies only the cell order, hence only the allocation $(t_c)$, which ranges
over prefix-filling allocations; relaxing to all allocations satisfying the
stated constraints can only increase the maximum.  That relaxation is a
fractional knapsack with unit weights, whose optimum sorts cells by $r_c$
descending---and this optimum is itself prefix-filling, so the relaxation is
tight.
\end{proof}

\begin{corollary}
\label{cor:nomethod}
No choice of minimisation algorithm, labelling threshold, don't-care policy,
cover-selection rule, or Boolean representation can raise the expected
precision of a predictor above the bound of Theorem~\ref{thm:ceiling} for its
feature map.  Improvement requires refining $\varphi$.
\end{corollary}

\begin{proof}
Each listed choice determines which cell predictor $\rho$ is obtained; none
changes $\varphi$.  The bound is uniform over $\rho$.
\end{proof}

Table~\ref{tab:ceiling} evaluates the bound on all 150 PSICOV targets at
$k = L/5$, with cells scored by their own pooled empirical contact rate---the
feature class handed the test labels.

\begin{table}[htbp]
\centering
\caption{Achievable prec@$L/5$ per feature map, 150 PSICOV targets.
Consensus residue codes; separation bins $[6,11)$, $[11,21)$, $[21,31)$,
$[31,\infty)$; mutual-information quartile computed per protein.}
\label{tab:ceiling}
\begin{tabular}{lrr}
\toprule
Feature map $\varphi$ & Cells & Achievable \\
\midrule
residue$_i$, residue$_j$, separation, MI quartile & 6390 & 0.2412 \\
residue$_i$, residue$_j$, separation              & 1600 & \textbf{0.1933} \\
physicochemical group$_i$, group$_j$, separation  & 256  & 0.1418 \\
residue$_i$, residue$_j$                          & 400  & 0.1383 \\
Gray Hamming distance, separation                 & 24   & 0.0789 \\
separation only                                   & 4    & 0.0598 \\
\midrule
\multicolumn{3}{l}{\emph{Measured on the same benchmark:}} \\
Trained circuit, held out (Sec.~\ref{sec:res-transfer}) & & 0.179 \\
Plain mutual information                                & & 0.097 \\
Published DCA~\cite{jones2012}                          & & 0.723 \\
\bottomrule
\end{tabular}
\end{table}

Three readings follow.  \textbf{The machinery is saturated}: the trained
circuit reaches $0.179$ against an achievable $0.1933$, or $93\%$ of its
feature class's ceiling, so by Corollary~\ref{cor:nomethod} the total remaining
headroom from every possible improvement in Boolean method is at most about
$8\%$ relative.  \textbf{The shortfall is informational}: the ceiling itself
sits $3.7\times$ below published DCA, and no Boolean construction on these
features can close that gap.  \textbf{Coarsening chemistry costs information}:
the eight-group physicochemical partition has a \emph{lower} ceiling than raw
residue identity, as it must, since merging cells can only reduce the bound.
Chemical classes buy interpretability, not accuracy.

This reframes Section~\ref{sec:res-reconstruction-negative}.  The failure to
meet the reconstruction criterion is not a deficiency of our minimiser or our
labelling; it is a property of consensus residue identity as an input.  Since
the bound depends only on the partition, the only lever is a finer partition
separating contacting from non-contacting pairs \emph{within} existing cells.
Adding a single crudely binned column statistic already lifts the ceiling from
$0.1933$ to $0.2412$, whereas richer residue vocabularies lift nothing---which
is the formal counterpart of the rank-fusion result of
Section~\ref{sec:res-gray-fusion}, where combining the circuit with continuous
DCA outperforms either alone.
